\section{Finite jumps and pointwise-small continuous singularities}
\label{sec:v44-jump-cusp}
We apply the complete height bound to the original density, before any
localization. It separates two different effects. A genuine jump keeps
its height when its support is shortened. A positive-exponent cusp
vanishes at its endpoint and can be localized with a prescribed
supremum budget, without scaling its coefficient. This distinction
allows an actual pointwise estimate for the entire continuous singular
part of the count-compatible correction.

\subsection{The germs permitted by the physical source}
\begin{corollary}[No divergent physical density germ]
\label{cor:v44-bounded-germs}
Fix $n,R$ and a bounded finite-record subanalytic insertion. For every
actual label $\ell=(k,m)$, the convergent expansion
\eqref{eq:power-log-expansion} of the complete component $f_\ell$ has
no nonzero term with negative exponent and no term
$(\log x)^j$ with $j\ge1$ and exponent zero. Its one-sided limits are
finite. The component, extended by zero, belongs to $\mathrm{BV}(\R)$.
The same exclusion holds at singular and decision-boundary values, not
only at regular critical values.
\end{corollary}
\begin{proof}
Theorem~\ref{thm:v44-complete-density-bound} bounds the absolute
component density. Its constructible representative is analytic on each
of the finitely many regular roof intervals, so an almost-everywhere
bound holds everywhere on each such interval. In a convergent
power--logarithm expansion, the smallest nonzero exponent and then the
largest logarithmic power at that exponent determine its magnitude
near zero. A negative exponent or a nonconstant zero-exponent
logarithmic polynomial would contradict boundedness. Thus near each
one-sided endpoint
\begin{equation}\label{eq:v44-bounded-expansion}
 f_\ell(s\pm x)=l_{s,\pm}
       +\sum_{0<\alpha\le1}\sum_j c_{s,\pm,\alpha,j}
                                      x^\alpha(\log x)^j+r_{s,\pm}(x),
\end{equation}
where the sum is finite and the remainder has the derivative bounds
\eqref{eq:power-log-remainder} with exponent greater than one.
Each positive-exponent term has an integrable first derivative. The
remainder does too. On closed subintervals away from the finitely many
endpoints the density is analytic, with bounded derivative. Adding the
finitely many finite jumps proves bounded variation. Apply the argument
to real and imaginary parts. It concerns the full component; separate
branch singularities must first be combined as in the finite-graph lemma.
\end{proof}

\begin{proposition}[Uniform finite-count variation and smoothing]
\label{prop:v44-finite-count-variation}
For each $m$ there is $V_m<\infty$, independent of $R$, such that
\begin{equation}\label{eq:v44-BV-budget}
 \sum_{n=1}^m\sum_k\|D_t p_{n,R}(k,m,\cdot)\|_{\TV}\le V_m.
\end{equation}
Here $D_t$ is a distributional derivative, not the covariance $D_R$.
The same assertion holds for any specified jointly globally subanalytic
finite-record family of weights with a common bound and a fixed finite
parameter description. No such uniform variation conclusion is asserted
for all merely measurable bounded weights.

For the even kernel $K_B=B K_1(B\cdot)$ of
Theorem~\ref{thm:v43-arithmetic-raw-criterion},
\begin{equation}\label{eq:v44-L1-unsmoothing}
 \sum_{n,k}\|p_{n,R}(k,m,\cdot)-K_B*p_{n,R}(k,m,\cdot)\|_1
       \le \frac{V_m}{B}\int_\R |s|\,|K_1(s)|\,ds.
\end{equation}
No growth rate for $V_m$ is asserted.
\end{proposition}
\begin{proof}
For fixed $m$, allow the finitely many rectangle endpoints defining the
section to be independent bounded parameters. In the fixed rational
charts their membership tests are semialgebraic with these endpoints
as coefficients, retaining the fixed factor $c^{-1}$ rather than
renormalizing the enlarged sections; finite chart cases also cover endpoints passing a
chart cut. Thus the full finite graph and cumulative component integral
in \eqref{eq:constructible-cumulative} are jointly constructible in
these parameters, $R$ and $t$, by \cite[Theorem 1.3]{CMInt}.
Differentiation in $t$ on a definable partition gives a definable density
representative in the o-minimal real field with restricted analytic
functions and exponentiation. The uniform monotonicity consequence of
cell decomposition \cite{Coste} bounds the number of monotonicity
intervals of every real and imaginary component by a finite number
depending on $m$ and the fixed description, not on the parameters.
The height bound is uniform. A function bounded by $U$ on $N$
monotonicity intervals, extended by zero, has variation at most
$4(N+1)U$, including endpoint jumps. There are only finitely many labels
and return indices at fixed $m$. This proves \eqref{eq:v44-BV-budget}.
The actual moving endpoints are a subfamily of this enlarged parameter
family; no unproved definability in $R$ of their chosen parametrization
is needed. A specified jointly subanalytic weight is treated by the
same parameter integral. This proof is not an effective long-time
bound on the monotonicity count.

For a compactly supported BV function $f$,
$\|f(\cdot-s)-f\|_1\le |s|\|Df\|_{\TV}$, by integration of its
derivative measure. Use $\int K_B=1$, the triangle inequality and
Tonelli, then sum \eqref{eq:v44-BV-budget}. Scaling the kernel gives
\eqref{eq:v44-L1-unsmoothing}. The kernel need not be positive.
\end{proof}

\subsection{A count-compatible decomposition with a supremum budget}
Fix $n,R,w$ as in Corollary~\ref{cor:v44-bounded-germs}, and put
$b_\ell=\nu_R^*\{(K_n,N_n)=(k,m)\}$. Null components are set to zero.
At a nonnull component let
$J_{\ell,s}=f_\ell(s+)-f_\ell(s-)$ be each intrinsic jump.

\begin{theorem}[Jump--cusp separation of the complete source]
\label{thm:v44-jump-cusp-extraction}
For every $\varepsilon,\delta>0$ there is an exact decomposition
\begin{equation}\label{eq:v44-jump-cusp-identity}
 f_\ell=j_\ell+z_\ell+q_\ell,
       \qquad q_\ell\in W^{2,1}(\R),
\end{equation}
with compactly supported components, such that
\begin{align}
 \|j_\ell\|_1+\|z_\ell\|_1&\le\varepsilon b_\ell,
                                            \label{eq:v44-edge-mass}\\
 \|z_\ell\|_\infty&\le\delta b_\ell,
 \qquad \|j_\ell\|_\infty\le\max_s|J_{\ell,s}|.
                                            \label{eq:v44-cusp-height}
\end{align}
The empty maximum is zero. The function $j_\ell$ consists only of
localized step functions with coefficients $J_{\ell,s}$.
The continuous function $z_\ell$ contains precisely the localized
positive-exponent terms of \eqref{eq:v44-bounded-expansion} through
exponent one. All intrinsic coefficients are unchanged.

The correction series converge absolutely in total variation and
\begin{equation}\label{eq:v44-summed-cusp}
 \sum_\ell\|z_\ell\|_\infty\le\delta.
\end{equation}
All collision-count truncations are restrictions of these same
components. Their total-variation tails and exponential count moments
satisfy the estimates of Theorem~\ref{thm:v43-complete-correction}.
\end{theorem}
\begin{proof}
Fix the complete component before choosing any neighborhoods. Choose a
$C^2$ cutoff $\theta$ on $[0,\infty)$, equal to one near zero,
zero beyond one, taking values in $[0,1]$, with matching derivatives
at the transition endpoints. At every singular value $s$ put
\[
 j_{\ell,s}(t)=J_{\ell,s}\one_{t>s}\theta((t-s)/d_s).
\]
Use disjoint two-sided neighborhoods of the distinct singular values.
Subtract on each side the positive-exponent jet in
\eqref{eq:v44-bounded-expansion}, multiplied by the corresponding
$\theta(x/d_s)$, to define $z_{\ell,s}$.
Every such jet tends to zero at its endpoint. Its supremum on
$0<x<d_s$ and its integral there tend to zero as $d_s\downarrow0$;
the jump term's integral also tends to zero. Shrink the finitely many
radii to make the sum of the integrals at most
$\varepsilon b_\ell$ and the sum of cusp suprema at most
$\delta b_\ell$. This changes no coefficient. Disjoint jump supports
and $0\le\theta\le1$ give the second bound in
\eqref{eq:v44-cusp-height}.

Near $s$, the residual on both sides has the same limiting value
$l_{s,-}$, because the right jump $l_{s,+}-l_{s,-}$ was removed.
Both residual first derivatives tend to zero, and their second
derivatives are integrable by the remainder estimate with exponent
greater than one. The constant common trace need not be subtracted.
At the cutoff transitions the values and first derivatives match.
Thus there is no atom in the residual's first or second distributional
derivative, and $q_\ell\in W^{2,1}$. Unlike the earlier independent
one-sided subtraction, this construction has a possibly nonzero,
common value trace; zero traces are not needed for Fourier inversion.

Sum the budgets over labels using $\sum b_\ell=1$ and
$\|f_\ell\|_1\le\|w\|_\infty b_\ell$. This proves all claimed
series bounds. Choose $d_s\le1$; the supports then stay in
$[-1,m\tau_{\max}+1]$. Restricting by $m\le L$ changes no component
already chosen, so the count compatibility is exact. The cumulative
return tail applied to $\sum_{m>L,k}b_{(k,m)}$ proves the same
variation tails; the enlarged roof support changes only a constant
in the exponential-moment calculation. Selection need not be continuous
in the radius for any of these conclusions.
\end{proof}

\begin{corollary}[A pointwise bound for the full continuous singular part]
\label{cor:v44-cusp-local-scale}
Let $B>0$ be fixed and let $\mathscr D_{B,n,R}$ be the unchanged
full-source signed correction. With $h_\ell=\widehat q_\ell$ from
\eqref{eq:v44-jump-cusp-identity}, it has the exact decomposition
\begin{align}\label{eq:v44-refined-remainder}
 \mathscr D_{B,n,R}(\ell,t)
   ={}&j_\ell(t)-(K_B*j_\ell)(t)
       +\frac1{2\pi}\int e^{-ibt}(1-\chi(b/B))h_\ell(b)\,db
       +r_\ell(t),\\
 \sum_\ell\|r_\ell\|_\infty
   &\le(1+\|K_1\|_1)\delta.\notag
\end{align}
The inverse is absolutely convergent after taking the label coefficient.
For $\delta_n=(1+n)^{-P-3}$ and central labels $m\asymp n$,
\[
 m^2\sum_\ell\|r_\ell\|_\infty=O(n^{-P-1}),
\]
uniformly in $R$. This estimate includes every positive-exponent
singular and decision-boundary germ of the complete component.
\end{corollary}
\begin{proof}
Set $r_\ell=z_\ell-K_B*z_\ell$. The convolution inequality and
\eqref{eq:v44-summed-cusp} give its bound; $\|K_B\|_1=\|K_1\|_1$.
The remaining identity follows from $f_\ell-K_B*f_\ell$ and the
$W^{2,1}$ coefficient-first inversion already proved. The common
nonzero trace of the residual introduces no distributional atom.
Multiplication by the central $m^2=O(n^2)$ proves the stated rate.
\end{proof}

\begin{proposition}[The jump size visible to raw unsmoothing]
\label{prop:v44-jump-test}
At an intrinsic jump $s$ and for every fixed $B>0$,
\begin{equation}\label{eq:v44-jump-lower-bound}
 \operatorname*{ess\,sup}_{|t-s|<h}|\mathscr D_{B,n,R}(\ell,t)|
                    \ge\tfrac12|J_{\ell,s}|\qquad(h>0).
\end{equation}
Thus a central pointwise local law with a continuous roof main term
requires $m^2|J_{\ell,s}|\to0$ for jumps lying in the interior of the
central target intervals.
\end{proposition}
\begin{proof}
The convolution of the compactly supported $L^1$ density with the
Schwartz kernel is continuous. The two one-sided traces of
$\mathscr D_B=f_\ell-K_B*f_\ell$ consequently differ by $J_{\ell,s}$.
At least one trace has modulus at least half this difference. Continuity
on the adjoining regular intervals gives the essential-supremum bound.
Apply the same argument to the difference from a continuous main term.
\end{proof}

\paragraph{What remains on the local scale.}
The theorem gives a height estimate for the actual source, excludes
divergent germs, and pays every continuous singular term in a summed
supremum norm. It does not bound the jump part or the residual inverse
in \eqref{eq:v44-refined-remainder} by $o(m^{-2})$. The selected radii
can still make $\|q_\ell''\|_1$ large. Their supremum budget is not a
bound for those derivatives, and the finite-count variation theorem
gives no usable growth rate for $V_m$. The main term remains
$\mathcal L_{m,R}$, or $c\mathfrak a_Rg_{\Omega_R}$ at fixed radius.
No arithmetic residue, prescribed singleton, conditioning event or
fixed-band quantifier has been changed.
